\documentclass[DIV16,BCOR1cm,10pt,a4paper,fleqn,twoside]{scrreprt}         % for pdf output
%\documentclass[DIV16,BCOR1cm,10pt,a4paper,fleqn]{scrreprt}         % for pdf output
%\documentclass[DIV16,BCOR1cm,11pt,a4paper,fleqn]{report}         % for pdf output

% To allow automatic selection of the right graphics type ...
% preset \pdfoutput for older latex installation, it is allways definted for
% news ones
\ifx\pdfoutput\undefined
\gdef\pdfoutput{0}
\fi

\newif\ifpdfx
\ifnum\pdfoutput=0
% latex is called for dvi output
   \pdfxfalse
   \usepackage{graphicx}
\else
% pdflatex is called for pdf output
   \pdfxtrue
   \usepackage[pdftex]{graphicx}
   \usepackage[pdftex]{hyperref}
\fi

\usepackage{textcomp}

%\newcommand{\CDO}{{\bfseries\sffamily CDO\ }}
\newcommand{\CDO}{{\bfseries\sffamily CDO}}
\newcommand{\cdologo}{\includegraphics{logo/cdo_logo}}

\graphicspath{{figures/}}

% To define headers and footers
\usepackage{fancyhdr}
\pagestyle{fancy}

% Headers and footers personalization using the `fancyhdr' package
\fancyhf{} % Clear all fields

\renewcommand{\headrulewidth}{0.2mm}
\renewcommand{\footrulewidth}{0.2mm}

\renewcommand{\chaptermark}[1]{\markboth{#1}{}}
\renewcommand{\sectionmark}[1]{\markright{#1}}

%\renewcommand{\chaptermark}[1]{\markboth{#1}{}}
%\renewcommand{\sectionmark}[1]{\markleft{#1}}

\fancyhead[LO,RE]{\slshape \leftmark}
\fancyhead[LE,RO]{\slshape \rightmark}
\fancyfoot[LE,RO]{\Large\thepage}
%\fancyfoot[LO,RE]{\raisebox{-2.8mm}{\scalebox{0.17}{\cdologo}}}
\fancypagestyle{plain}{%
  \fancyhead{} % get rid of headers
  \renewcommand{\headrulewidth}{0pt}
}

%\setlength{\footnotesep}{0cm}
%\setlength{\footskip}{-2cm}
%\renewcommand{\footnoterule}{\rule{0cm}{0cm}}

\usepackage{exscale}
\usepackage{array,colortbl}    % color table

\usepackage{listings}
\usepackage{longtable}
\usepackage{color}
\definecolor{pcolor1}{rgb}{0.992, 0.980, 0.875}  % rgb: 253/250/223
\definecolor{pcolor2}{rgb}{1.000, 0.925, 0.700}  % rgb: 255/236/278
\definecolor{pcolor3}{rgb}{0.968, 0.756, 0.623}  % rgb: 247/193/159

%\usepackage{ae}               % fuer die "almost european" computer modern fonts
%\usepackage{url}              % Standard-Paket fuer WWW-Adressen

%\typearea{10}                 % Einen sinnvollen Satzspiegel aktivieren


%Usage:
%pdflatex cdo.tex
%pdflatex cdo.tex
%cat > cdo.ist << 'EOF'
%delim_0        "{\\idxdotfill} "
%headings_flag  1
%heading_prefix "{\\centerline {\\Large \\textbf{"
%heading_suffix "}}}"
%EOF
%makeindex -s cdo.ist cdo.idx
%pdflatex cdo
%thumbpdf cdo
%pdflatex cdo


\usepackage{thumbpdf}

%\usepackage{html}

\usepackage{makeidx}

%\ifpdf
%\usepackage[a4paper, colorlinks=true, pdfstartview=FitV, bookmarks=true, linkcolor=blue,
%            citecolor=blue, urlcolor=blue, latex2html=true]{hyperref}
%\fi

\usepackage{hyperref}
\hypersetup{pdftoolbar=true,
            pdfmenubar=true,
            pdfwindowui=true,   
%            pdffitwindow=true,
            pdfauthor={Uwe Schulzweida},
            pdftitle={CDO Climate Data Operators},
            pdfcreator={pdflatex + hyperref},
            pdfstartview=FitV,
%            pdfpagemode=FullScreen,
            a4paper,
            bookmarks=true,
            linkcolor=blue,
            citecolor=blue,
            urlcolor=blue,
            colorlinks=true}

\setlength{\parindent}{0em}
\setlength{\parskip}{1.5ex plus0.5ex minus0.5ex}
\extrarowheight1pt

\makeindex
%\newcommand{\ii}[1]{\textit{#1}}  \newcommand{\nn}[1]{#1n}
%\renewcommand{\dotfill}{\leaders\hbox to 5p1{\hss.\hss}\hfill}
%\newcommand{\idxdotfill}{5p1{\hss.\hss}\hfill}
\newcommand{\idxdotfill}{\ \dotfill \ }
%\def\idxdotfill{\leaders\hbox to.6em{\hss .\hss}\hskip 0pt plus 1fill}
%\MakeShortVerb{\@}

\renewcommand{\indexname}{Operator index}

\newenvironment{defalist}[1]
{\begin{list}{}
{\settowidth{\labelwidth}{#1\ \ }
\setlength{\itemsep}{0mm}
\setlength{\itemindent}{0mm}
%\setlength{\listparindent}{25mm}
\setlength{\leftmargin}{\labelwidth}
%\setlength{\leftmargin}{25mm}
\setlength{\labelsep}{2mm}
\addtolength{\leftmargin}{\labelsep}
%\addtolength{\leftmargin}{8mm}
}}
{\end{list}}

\newenvironment{defalist2}[1]
{\begin{list}{}
{\settowidth{\labelwidth}{#1\ \ }
\setlength{\itemsep}{0mm}
\setlength{\itemindent}{0mm}
%\setlength{\listparindent}{25mm}
\setlength{\leftmargin}{\labelwidth}
%\setlength{\leftmargin}{25mm}
\setlength{\labelsep}{2mm}
\addtolength{\leftmargin}{\labelsep}
\addtolength{\leftmargin}{8mm}
}}
{\end{list}}

\newcommand{\miniwidth}{\textwidth}

\setcounter{secnumdepth}{3}

\begin{document}


\begin{titlepage}
\vspace*{50mm}
{\Huge\textbf{CMOR support in \CDO}}

\setlength{\unitlength}{1cm}
\begin{picture}(16,0.4)
\linethickness{1.5mm}
%\put(0,0.1){\line(1,0){15.85}}
\put(0,0.1){\line(1,0){16.3}}
\end{picture}
\begin{flushright}
\large\textbf{CMORizing of climate model data \\ February 2017}
\end{flushright}

\vfill

\Large\textbf{Karl--Hermann Wieners, Uwe Schulzweida}

\Large\textsl{Max Planck Institute for Meteorology}

\Large\textbf{Mathis Rosenhauer}

\Large\textsl{Deutsches Klimarechenzentrum (DKRZ)}

\begin{picture}(16,1)
\linethickness{1.0mm}
%\put(0,0.7){\line(1,0){15.85}}
\put(0,0.7){\line(1,0){16.3}}
\end{picture}
\end{titlepage}

\tableofcontents

\chapter{Introduction}

The Climate Data Operators ({\CDO}) software is a collection of operators
for standard processing of climate and forecast model data.

This document describes the additional {\CDO} operator cmor. This
operator is an interface to the CMOR library version 3 from PCMDI.
The CMOR support is available with {\CDO} release 1.8.0 or later.

The CMOR (Climate Model Output Rewriter) library comprises a set of
functions, that can be used to produce CF-compliant NetCDF files that 
fulfill the requirements of many of the climate community's standard
model experiments. These experiments are collectively referred to as
MIP's and include, for example, AMIP, CMIP, CFMIP, PMIP, APE, and IPCC 
scenario runs. The output resulting from CMOR is "self-describing" and
facilitates analysis of results across models.

Much of the metadata written to the output files is defined in
MIP-specific tables, typically made available from each MIP's web
site. CMOR relies on these tables to provide much of the metadata 
that is needed in the MIP context, thereby reducing the programming 
effort required of the individual MIP contributors.

The CDO operator cmor reads the data with the internal IO library CDI
and writes the NetCDF result directly with the CMOR library. 

\chapter{Building CDO with CMOR}

This section describes how to build and install {\CDO} with CMOR
support on a UNIX system.

\section{CMOR}
 
CMOR version 3 needs to be installed before building {\CDO}.
CMOR depends on the following external libraries:
netCDF-4.0 , HDF5, UDUNITS-2, zlib and uuid.
Make sure that you use exactly the same libraries for the {\CDO}
installation, otherwise the operator cmor will possibly not working correctly.

\section{Compilation}

First go to the {\CDO}  \href{https://code.zmaw.de/projects/cdo}{\tt download} page
({\tt https://code.zmaw.de/projects/cdo}) to get the latest distribution,
if you do not have it yet.
Compilation is done by performing the following steps:

\begin{enumerate}
\item Unpack the archive, if you haven't done that yet:
   
\begin{verbatim}
    gunzip cdo-$VERSION.tar.gz    # uncompress the archive
    tar xf cdo-$VERSION.tar       # unpack it
    cd cdo-$VERSION
\end{verbatim}

\item Configure {\CDO} with CMOR support:

The configuration depends on the location of the external libraries. Here is one example:

\begin{verbatim}
./configure --with-cmor=<CMOR root directory> --with-netcdf=<NetCDFroot directory> \
            --with-uuid --with-udunits2 LIBS=-lossp-uuid                           \
            CPPFLAGS="-I<CMOR root dir>/include/cdTime -I<CMOR root dir>/include/json-c"
\end{verbatim}

For an overview of other configuration options use

\begin{verbatim}
    ./configure --help
\end{verbatim}

\item Compile the program by running make:

\begin{verbatim}
    make
\end{verbatim}

\end{enumerate}

The program should compile without problems and the binary ({\tt cdo}) 
should be available in the {\tt src} directory of the distribution.


\section{Installation}

After the compilation of the source code do a {\tt make install},
possibly as root if the destination permissions require that.

\begin{verbatim}
    make install
\end{verbatim} 

The binary is installed into the directory {\tt $<$prefix$>$/bin}.
{\tt $<$prefix$>$} defaults to {\tt /usr/local} but can be changed with 
the {\tt --prefix} option of the configure script. 


%\chapter{\label{refman}CMOR operator reference manual}

This section gives a description of all {\CDO} operators to generate plots.
Related operators are grouped to modules.
For easier description all single input files are named \texttt{ifile} or \texttt{ifile1}, \texttt{ifile2}, etc.,
and an arbitrary number of input files are named \texttt{ifiles}.
All output files are named \texttt{ofile} or \texttt{ofile1}, \texttt{ofile2}, etc.


\hspace{3mm}

\input{ref_list_cmor}
\input{ref_man_cmor}

\chapter{\label{refman}CMOR operator reference manual}

\hspace{3mm}

%\input{ref_list_cmor}
\input{ref_man_cmor_mod}

\begin{thebibliography}{xx}


\bibitem[CDI]{CDI} \ \\
  \href{https://code.zmaw.de/projects/cdi}
       {Climate Data Interface},
  from the
  \href{http://www.mpimet.mpg.de}
       {Max Planck Institute for Meteorologie}


\bibitem[CDO]{CDO} \ \\
  \href{https://code.zmaw.de/projects/cdo}
       {Climate Data Operators},
  from the
  \href{http://www.mpimet.mpg.de}
       {Max Planck Institute for Meteorologie}


\bibitem[CMOR]{CMOR} \ \\
  \href{git://github.com/PCMDI/cmor.git}
       {Climate Model Output Rewriter},
  from the
  \href{https://www-pcmdi.llnl.gov}
       { Program For Climate Model Diagnosis and Intercomparison (PCMDI)}


\end{thebibliography}


In this simple example a prepared ECHAM NetCDF file with the two 3D
variables \textbf{ta} and \textbf{ua} are cmorized with the {\CDO} operator \textbf{cmor}. 
Here is the NetCDF header of the input file named example.nc:
\begin{lstlisting}[frame=single, backgroundcolor=\color{pcolor1}, basicstyle=\footnotesize]
dimensions:
   lon = 192 ;
   bnds = 2 ;
   lat = 96 ;
   plev = 17 ;
   time = UNLIMITED ; // (120 currently)
variables:
   float lon(lon) ;
      lon:standard_name = "longitude" ;
      lon:long_name = "longitude" ;
      lon:units = "degrees_east" ;
      lon:bounds = "lon_bnds" ;
   float lon_bnds(lon, bnds) ;
   float lat(lat) ;
      lat:standard_name = "latitude" ;
      lat:long_name = "latitude" ;
      lat:units = "degrees_north" ;
      lat:bounds = "lat_bnds" ;
   float lat_bnds(lat, bnds) ;
   double plev(plev) ;
      plev:standard_name = "air_pressure" ;
      plev:long_name = "pressure" ;
      plev:units = "Pa" ;
      plev:positive = "down" ;
   double time(time) ;
      time:standard_name = "time" ;
      time:bounds = "time_bnds" ;
      time:units = "days since 1850-1-1 00:00:00" ;
      time:calendar = "proleptic_gregorian" ;
   double time_bnds(time, bnds) ;
   float ta(time, plev, lat, lon) ;
      ta:units = "K" ;
   float ua(time, plev, lat, lon) ;
      ua:units = "m s-1" ;
\end{lstlisting}
The example file contains no global attributes and only the units
attribute for the variables \textbf{ta} and \textbf{ua}.

The {\CDO} \textbf{cmor} command is:
\begin{lstlisting}[frame=single, backgroundcolor=\color{pcolor1}, basicstyle=\footnotesize]
   cdo cmor,CMIP6_Amon.json,__info=cmor.rc example.nc
\end{lstlisting}
The follwing two output files will be produced:
\begin{lstlisting}[frame=single, backgroundcolor=\color{pcolor1}, basicstyle=\footnotesize]
CMIP6/ta_Amon_piControl_MPI-ESM_r1i1p1f1_185001-185912.nc
CMIP6/ua_Amon_piControl_MPI-ESM_r1i1p1f1_185001-185912.nc
\end{lstlisting}
CMIP6\_Amon.json is the CMIP6 table file for monthly mean atmosphere
data from PCMDI. The option \texttt{\_\_info=cmor.rc} will read settting from
the config file cmor.rc.
Here is the contents of this config file:
\begin{lstlisting}[frame=single, backgroundcolor=\color{pcolor1},  basicstyle=\footnotesize]
__inpath                 = "/<XXX>/CMOR/cmip6-cmor-tables/Tables"
__grid_table             = "CMIP6_grids.json"
_control_vocabulary_file = "CMIP6_CV.json"
_cmip6_option            = "CMIP6"

activity_id              = "CMIP"
outpath                  = "CMIP6"
experiment_id            = "piControl"
realization_index        = 1
initialization_index     = 1
physics_index            = 1
forcing_index            = 1
source_type              = "AOGCM"
sub_experiment           = "none"
sub_experiment_id        = "none"
mip_era                  = "CMIP6"
contact                  = "cmip5-mpi-esm@dkrz.de"
institution_id           = "MPI-M"
model_id                 = "MPI-ESM-MR"
references               = "ECHAM6: n/a; JSBACH: n/a; MPIOM: n/a;  HAMOCC: n/a;"
branch_time              = 0
branch_method            = "standard"

grid                     = "native atmosphere T63 gaussian grid (64x128 latxlon)"
grid_label               = "gn"
grid_resolution          = "5 km"

institution              = "Max Planck Institute for Meteorology"

run_variant              = "MPI-ESM-MR model output prepared for CMIP5 pre-industrial control"
source_id                = "MPI-ESM"
source                   = "MPI-ESM:"
output_file_template     = "<variable_id><table><experiment_id><source_id><variant_label>"

license                  = "N/A"

# Rename axis
__rename_plev            = plev17
\end{lstlisting}

NetCDF header of the output file ta\_Amon\_piControl\_MPI-ESM\_r1i1p1f1\_185001-185912.nc:
\begin{lstlisting}[frame=single, backgroundcolor=\color{pcolor1}, basicstyle=\footnotesize]
dimensions:
   time = UNLIMITED ; // (120 currently)
   plev = 17 ;
   lat = 96 ;
   lon = 192 ;
   bnds = 2 ;
variables:
   double time(time) ;
      time:bounds = "time_bnds" ;
      time:units = "days since 1850-1-1 00:00:00" ;
      time:calendar = "proleptic_gregorian" ;
      time:axis = "T" ;
      time:long_name = "time" ;
      time:standard_name = "time" ;
   double time_bnds(time, bnds) ;
   double plev(plev) ;
      plev:units = "Pa" ;
      plev:axis = "Z" ;
      plev:positive = "down" ;
      plev:long_name = "pressure" ;
      plev:standard_name = "air_pressure" ;
   double lat(lat) ;
      lat:bounds = "lat_bnds" ;
      lat:units = "degrees_north" ;
      lat:axis = "Y" ;
      lat:long_name = "latitude" ;
      lat:standard_name = "latitude" ;
   double lat_bnds(lat, bnds) ;
   double lon(lon) ;
      lon:bounds = "lon_bnds" ;
      lon:units = "degrees_east" ;
      lon:axis = "X" ;
      lon:long_name = "longitude" ;
      lon:standard_name = "longitude" ;
   double lon_bnds(lon, bnds) ;
   float ta(time, plev, lat, lon) ;
      ta:standard_name = "air_temperature" ;
      ta:long_name = "Air Temperature" ;
      ta:comment = "Air Temperature" ;
      ta:units = "K" ;
      ta:history = "2017-01-31T11:23:09Z altered by CMOR: replaced missing value flag (-9e+33) with standard missing value (1e+20)." ;
      ta:missing_values = 1.e+20f ;
      ta:_FillValue = 1.e+20f ;

// global attributes:
      :Conventions = "CF-1.6 CMIP-6.0" ;
      :activity_id = "CMIP" ;
      :additional_source_type = "CHEM AER" ;
      :branch_method = "standard" ;
      :branch_time = 0. ;
      :contact = "cmip5-mpi-esm@dkrz.de" ;
      :creation_date = "2017-01-31T11:23:09Z" ;
      :data_specs_version = "01.beta.36" ;
      :experiment = "pre-Industrial control" ;
      :experiment_id = "piControl" ;
      :forcing_index = 1 ;
      :frequency = "mon" ;
      :further_info_url = "http://furtherinfo.es-doc.org/CMIP6/MPI-M.MPI-ESM.piControl.none.r1i1p1f1" ;
      :grid = "native atmosphere T63 gaussian grid (64x128 latxlon)" ;
      :grid_label = "gn" ;
      :grid_resolution = "5 km" ;
      :history = "2017-01-31T11:23:09Z CMOR rewrote data to be consistent with CF standards and CMIP6 requirements." ;
      :initialization_index = 1 ;
      :institution = "Max Planck Institute for Meteorology" ;
      :institution_id = "MPI-M" ;
      :mip_era = "CMIP6" ;
      :model_id = "MPI-ESM-MR" ;
      :physics_index = 1 ;
      :product = "output" ;
      :realization_index = 1 ;
      :realm = "atmos atmosChem" ;
      :references = "ECHAM6: n/a; JSBACH: n/a; MPIOM: n/a;  HAMOCC: n/a;" ;
      :run_variant = "MPI-ESM-MR model output prepared for CMIP5 pre-industrial control" ;
      :source = "MPI-ESM:" ;
      :source_id = "MPI-ESM" ;
      :source_type = "AOGCM" ;
      :sub_experiment = "none" ;
      :sub_experiment_id = "none" ;
      :table_id = "Amon" ;
      :table_info = "Creation Date:(11 October 2016) MD5:77b461da88bd079bd1932066c0c7ad4a" ;
      :title = "MPI-ESM model output prepared for CMIP6" ;
      :variable_id = "ta" ;
      :variant_label = "r1i1p1f1" ;
      :license = "N/A" ;
      :cmor_version = "3.1.0" ;
      :tracking_id = "57138bf0-d93b-495e-82d0-cf287a0849c1" ;
\end{lstlisting}


\clearpage
\ifpdfx
\phantomsection
\addcontentsline{toc}{chapter}{\indexname}
\printindex
\else
\input{catalog}
\input{alphabetic_list}
\fi
\end{document}
